\section{What Is Sent: Representing and Encoding Immersive Content}\label{sec:encode}

Every networked XR system must first decide what its content is: a sphere of pixels, a set of camera
views, a cloud of points, the parameters of a learned scene model or a stream of force samples. The
design question is which representation travels and how it is compressed. The choice fixes four
quantities that every later function inherits. The \emph{bitrate floor},
below which quality collapses, sets what the network must supply (\cref{sec:provision}). The
\emph{encoder delay} counts against the interaction budgets of \cref{sec:compute,sec:connect}. The
\emph{decoder cost} lands on a battery-powered headset or phone. \emph{Viewport dependence} decides
whether the stream can be cut to what the user will see, and so whether prediction
(\cref{sec:anticipate}) has anything to act on.

We group representations into four families by what they send (\cref{tab:encode-families}).
\emph{Image-based} representations send what cameras see. \emph{Geometry-based} representations send
the shape of the scene. \emph{Learned} representations send a model of the scene, or its meaning, for
the receiver to regenerate. \emph{Signal-based} representations send a physical signal sampled by the
rules of the sense that receives it. \Cref{fig:encode-chains} follows one chain per family from capture
to display.

\begin{figure}[t]
\centering
\begin{adjustbox}{max width=\linewidth}\begin{tikzpicture}[font=\sffamily\scriptsize,
  st/.style={draw=encode!70!black, fill=encode!8, rounded corners=1.5pt, line width=0.5pt,
             text width=2.0cm, minimum height=0.8cm, align=center, inner sep=2pt},
  hv/.style={st, fill=encode!26, draw=encode!80!black, line width=1.1pt},
  flow/.style={-{Stealth[length=1.8mm]}, line width=0.55pt, draw=black!60},
  rate/.style={font=\sffamily\scriptsize, fill=encode!5, inner sep=1pt, align=center, text=black,
               text width=1.45cm},
  hd/.style={font=\sffamily\footnotesize\bfseries, align=center, text width=2.1cm},
  lane/.style={font=\sffamily\scriptsize\bfseries, anchor=south west, inner sep=0.5pt}]
  \def\xa{0} \def\xb{2.48} \def\xc{4.96} \def\xd{7.44} \def\xn{9.45} \def\xe{11.46} \def\xf{13.94}
  % network band
  \fill[encode!5] (8.6,0.62) rectangle (10.3,-6.62);
  \draw[densely dashed, encode!60!black, line width=0.5pt] (8.6,0.62) -- (8.6,-6.62);
  \draw[densely dashed, encode!60!black, line width=0.5pt] (10.3,0.62) -- (10.3,-6.62);
  \node[font=\sffamily\footnotesize\bfseries] at (\xn,0.25) {Network};
  % headers
  \foreach \x/\t in {\xa/Capture, \xb/Representation, \xc/Encoder, \xd/Packetise\\[-1pt]and transport,
                     \xe/Decoder, \xf/Render and\\[-1pt]display}
    \node[hd] at (\x,0.25) {\t};
  % one lane: id, y, name, capture, representation, encoder, transport, encoder style,
  % decoder, decoder style, renderer, renderer style, rate on the wire
  \newcommand{\lanerow}[8]{%
    \node[lane] at (-1.05,#2+0.43) {#3};
    \node[st] (c#1) at (\xa,#2) {#4};
    \node[st] (r#1) at (\xb,#2) {#5};
    \node[#8] (e#1) at (\xc,#2) {#6};
    \node[st] (t#1) at (\xd,#2) {#7};
    \draw[flow] (c#1) -- (r#1); \draw[flow] (r#1) -- (e#1); \draw[flow] (e#1) -- (t#1);
    \lanerest{#1}{#2}}
  \newcommand{\lanerest}[2]{\def\li{#1}\def\ly{#2}\lanerestb}
  \newcommand{\lanerestb}[5]{%
    \node[#2] (d\li) at (\xe,\ly) {#1};
    \node[#4] (p\li) at (\xf,\ly) {#3};
    \draw[flow] (t\li) -- (d\li); \draw[flow] (d\li) -- (p\li);
    \node[rate] at (\xn,\ly) {#5};}
  \lanerow{1}{-0.75}{Image-based: 360-degree video}{camera rig,\\stitching}{ERP or CMP,\\region packing}
    {HEVC/VVC tiles,\\quality ladder}{OMAF/DASH\\tile segments}{st}
    {decoder per tile\\(CPU 18$\to$58\%)}{hv}{viewport\\extraction}{st}{6--22 Mbps}
  \lanerow{2}{-2.0}{Image-based: multi-view plus depth (MIV)}{colour + depth\\cameras}{view atlases\\(MIV/V3C)}
    {HEVC or H.264\\per stream}{RTP or DASH,\\several streams}{st}
    {hardware video\\decoders}{st}{view synthesis\\or fusion}{hv}{30--100\\Mbps}
  \lanerow{3}{-3.25}{Geometry-based: point cloud}{RGB-D camera\\array}{points, octree,\\3D tiles}
    {V-PCC (slow),\\G-PCC, Draco}{DASH/WebRTC,\\per tile}{hv}
    {video decode,\\3D rebuild}{hv}{point splatting,\\6DoF view}{st}{19 Mbps;\\raw $\leq$5 Gbps}
  \lanerow{4}{-4.5}{Learned: Gaussian or neural scene}{multi-view\\video}{3DGS/NeRF or\\semantic tokens}
    {train each frame\\($\approx$1 min)}{baked into\\2D video}{hv}
    {H.264 or neural\\decoder}{st}{splat rendering\\(37 ms/frame)}{hv}{0.53 MB\\per frame}
  \lanerow{5}{-5.75}{Signal-based: haptic}{force sensor,\\$\geq$1 kHz}{kinaesthetic or\\vibrotactile}
    {perceptual\\deadband}{event-driven\\small packets}{st}
    {hold or predict\\lost samples}{st}{actuator in the\\control loop}{st}{packets\\$-$80--90\%}
  % sender / receiver braces
  \draw[decorate, decoration={brace, amplitude=3pt, mirror}, draw=black!55]
    (-1.05,-6.3) -- (8.5,-6.3) node[midway, below=3pt, font=\sffamily\scriptsize] {sender: capture and encoding};
  \draw[decorate, decoration={brace, amplitude=3pt, mirror}, draw=black!55]
    (10.4,-6.3) -- (15.0,-6.3) node[midway, below=3pt, font=\sffamily\scriptsize] {receiver: headset, phone or actuator};
  % legend
  \node[hv, text width=0.35cm, minimum height=0.28cm] (lg) at (-0.75,-7.05) {};
  \node[anchor=west, font=\sffamily\scriptsize] at (lg.east) {stage that dominates computation};
  \node[rate, text width=0.75cm, draw=encode!40] (lr) at (5.3,-7.05) {Mbps};
  \node[anchor=west, font=\sffamily\scriptsize] at (lr.east) {bitrate or packet rate on the wire (measured examples)};
\end{tikzpicture}\end{adjustbox}
\caption{Transmission chains from capture to display for one or two representations of each family. No
chain is cheap at both ends: a low rate on the wire is bought with encoder delay (V-PCC, per-frame
training) or with computation on the receiver (tile decoding, view synthesis, splat rendering). Figures:
360-degree video \cite{W2623027537,W2896840651}, multi-view streams \cite{W4200107448}, point clouds
\cite{W4362496249,W3021314748}, Gaussian video \cite{W4404525293}, deadband coding \cite{W2912555763}.}
\label{fig:encode-chains}
\end{figure}

\begin{table}[t]
\centering\small
\caption{Representation families and where their cost lands. Figures are examples measured on the cited
systems under different conditions; they locate each cost but are not matched operating points.}
\label{tab:encode-families}
\setlength{\tabcolsep}{3pt}
\renewcommand{\arraystretch}{1.08}
\begin{tabularx}{\linewidth}{@{}>{\raggedright\arraybackslash}p{1.6cm}>{\raggedright\arraybackslash}X%
>{\raggedright\arraybackslash}p{2.3cm}>{\raggedright\arraybackslash}p{2.4cm}%
>{\raggedright\arraybackslash}p{2.5cm}>{\raggedright\arraybackslash}p{2.55cm}@{}}
\toprule
Family & Idea & Cost on wire & Cost at sender & Cost at receiver & Maturity\\
\midrule
Image & 360-degree video: sphere projected to rectangles, bits on the viewport & 6--22~Mbps & low; tiles lose 7--10\% efficiency & a decoder per tile, CPU 18$\to$58\% & standard (OMAF), deployed\\
\addlinespace[2pt]
Image & Multi-view plus depth (MIV): camera views, new views synthesised & 30--100~Mbps & low; hardware video encoders & several decoders and view synthesis & standard; rarely tested on networks\\
\addlinespace[2pt]
Geometry & Point clouds (V-PCC, G-PCC, Draco) & 19~Mbps live; raw up to 5~Gbps & V-PCC hours per second of content; G-PCC, Draco fast & 3D rebuild; 30~fps on phones needs inter-frame coding & standard; on-demand prototypes\\
\addlinespace[2pt]
Geometry & Meshes: surfaces from fused depth & 14~Mbps mean (reconstructed scene) & reconstruction & rasterisation & static codecs benchmarked\\
\addlinespace[2pt]
Learned & Neural scene (NeRF, 3DGS): model parameters per frame & 0.53~MB per frame & training $\approx$1~min per frame & splat rendering, 37~ms per frame & research prototypes\\
\addlinespace[2pt]
Learned & Semantic and generative: meaning, not pixels & 93--99.9\% less & inference, shared knowledge base & generative inference; 13.7~s end to end & simulation, one 5G prototype\\
\addlinespace[2pt]
Signal & Haptic: samples sent only above a perceptual threshold & 80--90\% fewer packets & negligible, delay-free & hold or predict missing samples & standard (IEEE~1918.1.1)\\
\addlinespace[2pt]
Signal & Holographic: light field or phase pattern & 6.63~TBps by arithmetic; 173~Mbps measured & V-PCC 11--42~min per second & display not characterised & requirement tables\\
\bottomrule
\end{tabularx}
\end{table}

\subsection{Image-based: send what a camera sees}
Image-based representations keep the 2D video chain and its hardware codecs. Their open questions are
how the scene maps onto rectangles and how many decoders the client can afford.

\emph{Projected 360-degree video} assumes a viewer who turns the head but stays in place (3DoF), and wins
when that holds. Concentrating bits on the viewport saves bitrate, but less than it saves pixels:
Oculus's offset cube used fewer than half the pixels of an equirectangular frame, yet saved only
5.6--16.4\% of bitrate \cite{W2623027537}. The delivery unit is contested. Viewport-specific versions
multiply storage; tiles share one encode but cost 7.3--10.4\% in coding efficiency, and Flare's tile
decoders raised phone CPU load from 18\% to 58\% \cite{W2896840651}. Gaze is the stronger lever.
Gaze-based quantisation saved 40--47\% of bitrate in cloud VR gaming \cite{W3174397927}, and
gaze-contingent compression was five times more effective when motion-to-photon (MTP) latency stayed
below 15~ms \cite{W4214531162}. Foveation is thus a coding gain that \cref{sec:compute} pays for in
latency.

\emph{Multi-view plus depth} sends colour and depth streams and lets the receiver synthesise the view,
giving limited six degrees of freedom (6DoF) with ordinary video codecs. It assumes the receiver can
afford view synthesis, and it wins in live, two-way systems because it needs no slow 3D encoder. Project
Starline sends seven H.265 streams at 30--100~Mbit/s and fuses them at the receiver, because real-time
3D compression was impossible at its resolution \cite{W4200107448}. MPEG Immersive
Video (MIV) standardises this path \cite{W4210296418}, yet no study has streamed MIV to users over a real
network.

\subsection{Geometry-based: send the shape}
Geometry-based representations send the scene itself, so the receiver can render any viewpoint. The
price is weight: raw point clouds need 720~Mbit/s at 200K points and 30~fps \cite{W3021314748}, one to
three orders of magnitude more than 2D video.

Point-cloud codecs trade compression ratio against speed, and the trade is established. V-PCC projects
patches onto planes and reuses video codecs; it reaches lossy ratios of 100--500 but took about 4~hours
to encode 30 frames \cite{W3096788322}. G-PCC and Draco code geometry in 3D, run fast and compress less,
so live systems avoid V-PCC. In VR2Gather, the MPEG Anchor codec cut 281.8 to 19.0~Mbit/s but raised
latency from 40.6 to 278.7~ms, and parallel tiles brought it back to 90.6~ms \cite{W4362496249}: coding,
not transit, set the delay. The receiver is the second constraint: Hermes reached 30~fps on phones only
by exploiting inter-frame correlation \cite{W4387967921}. No codec wins everywhere; in a comparison of
five codecs the best choice depended on content and viewing behaviour \cite{W4321488303}.

Geometry's advantage is a richer viewport dependence: points can be culled by visibility, occlusion and
distance, not only by direction. ViVo saved 40\% of bandwidth at SSIM above 0.99, but only
10\% for phone users \cite{W3021314748}. Meshes are less explored: a reconstructed mesh for robot teleoperation needed
14~Mbit/s on average \cite{W2966611292}.

\subsection{Learned: send a model or a meaning}
Learned representations replace explicit content with something the receiver must regenerate. They
promise the largest savings and move the most work to the two ends.

\emph{Neural scene representations} (NeRF, 3D Gaussian splatting) send model parameters. The most
deployable design reuses existing hardware: V$^3$ bakes dynamic Gaussians into 2D video for H.264
decoders and reaches 32.97~dB PSNR at 0.532~MB per frame and 27~fps on a phone \cite{W4404525293}. At
30~fps that is about 128~Mbit/s, the range of multi-view video, so the gain is rendering freedom rather
than bitrate. Training takes about a minute per frame, which rules out live content, and error control and
quality assessment remain open \cite{W4408847683}.

\emph{Semantic and generative codecs} send meaning and assume that sender and receiver share a knowledge
base or a generative model. Their comparisons are the weakest of any family. WiserVR sent
1.81 instead of 28.3~Mbit for 20 frames, in an AWGN simulation that ignored the cost of sharing the
knowledge base \cite{W4324296656}. The one generative prototype timed end to end cut bandwidth by
99.93\% but raised latency from 980~ms to 13.66~s \cite{W4403677254}. Fidelity metrics differ between
studies and can disagree with perceived realism \cite{W4385452934}, and baselines are usually BPG with
LDPC rather than HEVC or VVC. An order-of-magnitude saving is plausible; its price in inference delay
and energy is assumed rather than measured.

\subsection{Signal-based: send the physical signal}
Signal-based representations code a physical quantity by the limits of the sense that receives it.
Haptic and holographic content sit at opposite ends of the rate scale.

\emph{Haptic signals} need little rate and strict timing. Kinaesthetic samples travel at 1~kHz or more to
keep the control loop stable, and perceptual deadband coding sends a sample only when its change exceeds
a Weber fraction (7--15\% for force), removing 80--90\% of packets without perceptible loss. Two
independent implementations of the IEEE~1918.1.1 codec reproduced the same rates and errors
\cite{W2912555763}, an established result, though for one contact point. Vibrotactile coding is
open-loop and tolerates tens of milliseconds of transform delay \cite{W2911218692}; learned tactile
codecs report 80--92\% cuts, but only against other learned codecs \cite{W7131889935}. Deadband coding
makes traffic bursty, unlike the periodic 1~ms flow that transport studies assume.

\emph{Holographic content} has no agreed representation, so its requirement figures depend on an unstated
one. A light-field display example needs 221~GB per frame, or 6.63~TBps at 30~fps
\cite{W4300960574}. Measured systems used 5--100~Mbps for compressed point clouds over live 5G
\cite{W4312662428} and 173~Mbps for light-field telepresence \cite{W3176472668}: one figure is
uncompressed display arithmetic, the other compressed capture. For live use, codec delay binds first:
V-PCC needs 11--42 minutes to encode one second \cite{W4309030584}.

\paragraph{The choice moves cost; it does not remove it.} No representation is cheap at both ends.
Image-based content is light on the sender and heavy on the wire and the decoders. Geometry cuts the wire
at the price of encoder delay. Learned representations cut it further and move the work into training
and inference. Choosing a representation is choosing where to spend: on the network, the
sender's encoder or the receiver's processor. Nor is the choice independent. Remote rendering
(\cref{sec:compute}) collapses any representation back into 2D video at the edge, trading decoder cost
for MTP latency; viewport dependence bounds what prediction can save (\cref{sec:anticipate}); encoder
delay decides which representations two-way telepresence can use (\cref{sec:connect}). The ratio--speed
trade within each family is established; that new representations win across families is assumed, since
none has been compared with the others at matched quality.

\begin{openproblems}
\item \textbf{Which representation minimises end-to-end cost?} Point clouds, meshes,
  multi-view video and Gaussian splats have never been compared on the same content at matched perceived
  quality. Streaming shared scenes over replayed 5G and Wi-Fi traces would measure rate, encoder
  delay and decoder load.
\item \textbf{Do semantic and generative savings survive latency?} Reported savings come from AWGN
  simulations with weak baselines. Comparing semantic codecs with HEVC, VVC and V-PCC at matched quality
  on a real 5G link, timing inference at both ends, would show whether they suit interactive use.
\item \textbf{How should new representations be judged?} No rating data exist for compressed Gaussian,
  neural or generative reconstructions viewed in a headset, so PSNR and LPIPS are unvalidated for their
  artefacts. A 6DoF rating study across codecs and rates would provide ground truth.
\item \textbf{Can codecs track a mobile link?} Rate control is tuned at fixed rates, and neural codecs lack
  bitrate ladders and loss handling. Rate control driven by measured capacity, judged by frame-deadline
  misses, would test whether codec gains hold on cellular links.
\item \textbf{What does a hologram actually cost to send?} Requirement tables rest on display arithmetic,
  not captured content. Coding one scene as point clouds, splats, light fields and phase holograms and
  recording the traffic would replace assumed figures with measured ones.
\item \textbf{How do coded haptic streams behave under real loss?} Haptic codecs are evaluated on one-off
  recordings with one contact point. Shared multi-point recordings and a user study of deadband and
  vibrotactile codecs over replayed impairments would test perceptual transparency.
\end{openproblems}
